% claim-ref: cherry-shrimp#direct-development
% claim-ref: cherry-shrimp#incubation-disagreement
% claim-ref: cherry-shrimp#newborn-size
% claim-ref: cherry-shrimp#maturity-uncertainty
% claim-ref: cherry-shrimp#breeding-sequence
% claim-ref: cherry-shrimp#fish-risk
% claim-ref: cherry-shrimp#kuhli-risk
% claim-ref: cherry-shrimp#adult-transfer-gate
% claim-ref: cherry-shrimp#no-release
% claim-ref: cherry-shrimp#transfer-size
\CompanionHeader{REPRODUCTION}{Cherry shrimp}{Neocaridina davidi / freshwater colony planning}{Miniature hatchlings; no free-swimming saline larval rearing stage. [CS-UF]}
\Context{Keep the proposed nursery shrimp-only. The 20-gallon tank's Kuhli loaches are a separate compatibility decision; healthy adult coexistence cannot establish juvenile survival.}
\Panel{1. Establish readiness before breeding}{Confirm cycle evidence, stable water, food surfaces, protected intakes and healthy adults of both sexes. A purchase count cannot guarantee sex ratio. Arrange space and humane rehoming before the colony expands. \textbf{[CS-FARM; CS-UF; proposed]}}
\Panel{2. Observe eggs without handling}{Females hold eggs under the abdomen. UF/IFAS reports 16--19 days from literature; Co-Op says about a month, depending on temperature. Keep these estimates separate and record dates; do not force a hatch schedule. \textbf{[CS-UF; CS-COOP]}}
\Panel{3. Support the tiny juveniles}{Hatchlings are about 1--2 mm. Preserve biofilm and offer modest fine food near cover; avoid excess powder. Record hatchlings seen, feeding and growth separately. UF/IFAS reports maturity around 30 days, but home timing is unverified; do not raise heat to hurry breeding. \textbf{[CS-UF; CS-FARM; proposed]}}
\Panel{4. Review any adult transfer to the 20-gallon}{Proposed: compare both tanks' water, total stock, footprint and health history; arrange acclimation for actual water/transport differences and a fallback home. Retain a nursery breeding group. No safe transfer size or count is established. The nursery is not also a fish quarantine. \textbf{[CS-FARM; CS-VET; CS-KUH]}}
\Panel{Kuhli caveat / cover reduces exposure, not all risk}{Keeper accounts are reassuring, but are not predation trials. Pangio micropredator feeding plus shrimp vulnerability justify caution for juveniles and fresh molts. Even peaceful fish may eat tiny shrimp; a planted community is not a guaranteed nursery. \textbf{[CS-KUH; CS-SF; CS-COOP; CS-FARM; inference]}}
\Panel{Keep surplus contained}{Plan responsible rehoming; never release shrimp or aquarium contents outdoors. Introduced Neocaridina populations are documented. \textbf{[CS-UF]}}
\Panel{Observation record}{Eggs seen / hatchlings seen / feeding: \hrulefill\par
Water / growth / next decision: \hrulefill}
\Sources{Sources (clickable): \href{https://ask.ifas.ufl.edu/publication/IN1301}{CS-UF: UF/IFAS species profile}; \href{https://www.aquariumcoop.com/blogs/aquarium/cherry-shrimp-care}{CS-COOP: Co-Op cherry care}; \href{https://www.theshrimpfarm.com/posts/neocaridina-shrimp-care-breeding/}{CS-FARM: Shrimp Farm care}; \href{https://www.msdvetmanual.com/exotic-and-laboratory-animals/aquarium-fish/management-of-aquarium-fish}{CS-VET: MSD veterinary care}; \href{https://www.aquariumcoop.com/blogs/aquarium/care-guide-for-kuhli-loaches}{CS-KUH: Co-Op Kuhli care}; \href{https://www.seriouslyfish.com/species/pangio-semicincta/}{CS-SF: Seriously Fish Pangio}. Scopes, derivations and unknowns: adjacent YAML worksheets.}{cherry-shrimp / reproduction}
